\documentclass[UTF8,12pt]{ctexart}
\usepackage[a4paper,margin=2.3cm]{geometry}
\usepackage{array}
\usepackage{longtable}
\usepackage{booktabs}
\usepackage{hyperref}

\title{第八届中国研究生人工智能创新大赛\项目文档}\
\author{[团队名称]}
\date{2026.09.01}

\begin{document}
\maketitle

\section{项目概述}

\subsection{背景与问题}
SoulDrive 面向本机 Windows 单用户场景，聚焦“本地知识问答、多模态交互、项目空间隔离与证据可追溯”四类核心需求。现有通用问答系统在企业/竞赛型本地知识库场景中常见问题包括：检索不稳定、跨项目串扰、证据无法定位、多模态输入与文本回答脱节，以及一旦生成后端环境异常便无法形成完整闭环。针对这些痛点，SoulDrive 将检索、证据门、生成和多模态交互统一到一条可评测、可审计、可回放的本地链路中。

\subsection{项目价值}
本项目的价值体现在三点：
\begin{itemize}
  \item 以项目空间为最小隔离单元，实现本地知识、会话、实验与产物的分域管理，减少跨项目污染；
  \item 以证据驱动的 RAG 链路代替“直接生成”，保证回答可定位、可复核、可拒答；
  \item 以语音、屏幕、OCR、图像等多模态输入增强任务理解，使交互不再局限于纯文本问答。
\end{itemize}

\subsection{支持条件}
项目已完成本地运行链路的主体实现，包括知识入库、分层检索、图谱检索、证据门控、回答生成、运行监控与多模态工作流。当前验证以本机 Windows 环境为主，可离线运行或接入本地 Ollama 进行生成，适合竞赛提交和现场演示。

\section{项目规划}

\subsection{总体目标}
本项目目标是在单机本地条件下，构建一个具备“高准确率 RAG + 强证据约束 + 多模态交互 + 可追溯评测”的智能助手系统。系统需能够：
\begin{itemize}
  \item 正确识别用户问题所属项目与时间范围；
  \item 对本地知识库执行混合检索、重排与证据汇总；
  \item 对证据不足的问题稳定拒答；
  \item 对屏幕、语音与视觉输入形成统一上下文；
  \item 输出可审计、可复核、带引用的答案。
\end{itemize}

\subsection{创新点}
\begin{enumerate}
  \item \textbf{项目级 RAG 路由}：先按项目、资产类型和时间语义过滤，再决定是否进入复杂检索，减少无关召回。
  \item \textbf{混合证据链}：同时使用关键词、向量和图谱召回，并通过证据门校验标题命中、细节命中和引用有效性。
  \item \textbf{多模态任务对齐}：将语音、选区、OCR 和视觉证据纳入统一事件图，避免“看见了但答偏了”。
  \item \textbf{本地优先与可恢复}：生成后端可在 GGUF 与 Ollama 间切换，缺依赖时仍可用抽取式证据回答兜底。
\end{enumerate}

\subsection{应用场景}
\begin{itemize}
  \item 本地知识库问答与企业资料检索；
  \item 项目空间中的实验、文档、会议和任务追踪；
  \item 多模态场景下的屏幕内容理解与辅助决策；
  \item 参赛演示中的证据式问答与可解释展示。
\end{itemize}

\section{实施情况}

\subsection{技术路线}
系统采用“解析入库—分层索引—查询路由—混合检索—重排与证据门—受控生成—评测回归”的流水线设计：
\begin{enumerate}
  \item 文档解析与结构化抽取，保留标题、章节、页码、来源、项目标识等元数据；
  \item 子段落/父文档式切片，建立关键词索引、向量索引和图谱索引；
  \item 按查询意图选择直接回答、改写、多查询、分解或图谱检索；
  \item 对候选证据进行重排与门控，只保留可支撑答案的内容；
  \item 生成阶段强制引用证据，证据不足则拒答或抽取式兜底；
  \item 以评测集和回归测试验证检索、拒答、引用与多模态链路稳定性。
\end{enumerate}

\subsection{实施细节}
当前项目已完成：
\begin{itemize}
  \item 项目级检索与全局回退机制；
  \item GraphRAG 路由与图谱上下文控制；
  \item 证据门、引用校验与答案兜底；
  \item 运行态监控与模型后端切换；
  \item 多模态工作流与局部测试隔离修复；
  \item 全量测试回归通过。
\end{itemize}

\subsection{进度与分工}
建议在最终提交版中补充团队成员分工表，例如：
\begin{longtable}{p{0.28\textwidth}p{0.62\textwidth}}
\toprule
成员 & 分工 \\
\midrule
[成员1] & RAG 链路、检索评测与证据门 \\
[成员2] & 多模态输入、UI 与交互链路 \\
[成员3] & 文档、演示、实验与提交材料 \\
\bottomrule
\end{longtable}

\section{参考说明}

\subsection{可复核结果}
截至当前版本，项目已通过全量测试回归，说明主要功能模块在本机环境下可稳定运行。若最终比赛提交需附演示视频、架构图、指标截图或样例问答，请在正式封版前完成一次同环境复测并保留结果文件。

\subsection{提交前还需你确认的事项}
\begin{itemize}
  \item 团队名称、项目名称、版本日期、作者信息；
  \item 是否采用“Word 先审，再 LaTeX 出版”的流程；
  \item 是否需要我继续补：摘要、项目简介、实验结果、系统截图、参考文献；
  \item 是否要我把这份草稿直接转成 Word 可编辑稿或正式 PDF。
\end{itemize}

\end{document}
